\section{Diagonal jets and contact transitivity}
\label{jet:section}

Corollary~\ref{des:family} provides a symmetric section
\[
B\in H^0(X\times X,L\boxtimes L)
\]
and the identity
\begin{equation}\label{jet:family}
B(h(z),h(w))
=(z-w)^2\theta(h'(z))\otimes\theta(h'(w))
\end{equation}
on the full neighborhood $\mathcal H$ of the very free parametrized map $g$ of
Proposition~\ref{sm:very-free}. Here $z,w$ belong to one affine
coordinate chart of $\PP^1$, and $h'$ denotes differentiation in that
coordinate. We first determine the diagonal symbol of $B$ at every
point of $X$. We then prove a criterion that applies to any compact
complex contact manifold with such a symmetric section.

\subsection{From the full family to the diagonal}

\begin{lemma}\label{jet:value-velocity}
Fix the parameter $0\in\PP^1$ and its affine coordinate. There is a
Zariski-open neighborhood $\mathcal U\subset\mathcal H$ of $g$ on which
the value--velocity morphism
\[
\mathcal U\longrightarrow T_X,
\qquad h\longmapsto\bigl(h(0),h'(0)\bigr),
\]
is smooth. Its image is open in $T_X$. Its projection is a nonempty
open subset $U\subset X$, and for each $x\in U$ the velocities of maps
in $\mathcal U$ with $h(0)=x$ form a nonempty open subset of $T_xX$.
\end{lemma}

\begin{proof}
Put $E=g^*T_X$. Every splitting summand of $E$ has degree at least
one, so $H^1(E(-2\cdot0))=0$. The exact sequence
\[
0\longrightarrow E(-2\cdot0)\longrightarrow E
\longrightarrow E|_{2\cdot0}\longrightarrow0
\]
therefore makes $H^0(E)\to H^0(E|_{2\cdot0})$ surjective.
The Hom space is smooth at $g$ and has tangent space $H^0(E)$.
In target coordinates its joint value and velocity differential is
\[
\psi\longmapsto\bigl(\psi(0),\psi'(0)\bigr),
\]
which is precisely this restriction map, expressed in the
coordinate trivialization along $g$. It is surjective. The source and
target are smooth at the points in question, so the smooth locus of this
morphism contains a Zariski-open neighborhood $\mathcal U\subset\mathcal H$
of $g$. Its image in $T_X$ is open.
Projection to $X$ and intersection with its fibers give the final
assertions.
\end{proof}

Let $\Delta\subset X\times X$ be the diagonal, let $\mathcal I$
be its ideal sheaf, and put $M=L\boxtimes L$. We identify
$\mathcal I/\mathcal I^2$ with $\Omega_X^1$ by sending the class of
$f(y)-f(x)$ to $df$. Since $\Delta$ is a regular embedding,
\begin{equation}\label{jet:symbol-bundles}
(\mathcal I^k/\mathcal I^{k+1})\otimes M|_\Delta
\simeq\Sym^k\Omega_X^1\otimes L^{\otimes2}.
\end{equation}
We write $\theta^2$ for the symmetric square of the $L$-valued
one-form $\theta$, with quadratic value $v\mapsto\theta(v)^2$.

\begin{proposition}\label{jet:diagonal}
The section $B$ vanishes to order at least two along $\Delta$, and
its second diagonal symbol is
\begin{equation}\label{jet:second-symbol}
B\in H^0(X\times X,\mathcal I^2\otimes M),
\qquad [B]_2=\theta^2
\quad\text{in }H^0(X,\Sym^2\Omega_X^1\otimes L^{\otimes2}).
\end{equation}
\end{proposition}

\begin{proof}
Choose local coordinates and one local frame $e$ of $L$, used in
both factors. Write $\theta=e\alpha$ and
$B(x,y)=b(x,y)e(x)\otimes e(y)$. Fix $x\in U$ as in
Lemma~\ref{jet:value-velocity}. Setting $z=0$ in
Equation~\eqref{jet:family}, for a nearby map with
$h(0)=x$ and $h'(0)=a$, gives
\begin{equation}\label{jet:curve-expansion}
b(x,h(w))
=w^2\alpha_x(a)\alpha_{h(w)}(h'(w))
=w^2\alpha_x(a)^2+O(w^3).
\end{equation}
The constant term proves $B|_\Delta=0$ on $U$. As this restriction
is a global holomorphic section of $L^{\otimes2}$ on connected
$X$, it vanishes everywhere.

The first symbol is now intrinsically a section of
$\Omega_X^1\otimes L^{\otimes2}$. On $U$, its value on $a$ is
the coefficient of $w$ in Equation~\eqref{jet:curve-expansion}.
That coefficient vanishes for a nonempty open set of velocities,
so the linear form is zero. The first symbol consequently
vanishes on all of $X$, again by the identity Theorem. Thus
$B$ belongs to $\mathcal I^2\otimes M$ globally.

Its second symbol is now the global section in
Equation~\eqref{jet:symbol-bundles} for $k=2$. Because the linear
term has vanished, the acceleration $h''(0)$ contributes nothing
to the coefficient of $w^2$ in Equation~\eqref{jet:curve-expansion}.
The quadratic symbol takes the value $\alpha_x(a)^2$ on a nonempty
open set of $a$, and hence equals $\alpha_x^2$ as a quadratic form.
The difference $[B]_2-\theta^2$ is a holomorphic section of the
displayed symbol bundle, zero on $U$; it is therefore zero on $X$.
Each continuation has been performed after the preceding lower
symbol vanished, so all three identities are intrinsic.
\end{proof}

\subsection{A symmetric-jet criterion}

The remainder of the argument uses only symmetry and
Equation~\eqref{jet:second-symbol}. The quadratic symbol records contact
values of tangent vectors. Symmetry determines the next Taylor term.
Differentiating this expansion in the first slot reveals the Levi form,
which will detect directions inside the contact plane. We first make this local
calculation, then use it to prove Theorem~\ref{jet:criterion}.

\begin{lemma}\label{jet:local-calculation}
Let $B$ be symmetric and satisfy Equation~\eqref{jet:second-symbol}
on a complex contact manifold. In local coordinates and a frame
as above, put $s=y-x$ and
$A_x(u,v)=(\partial_u\alpha)_x(v)$, where $\partial$ differentiates
the coefficients of $\alpha$. Then
\begin{equation}\label{jet:cubic}
b(x,x+s)=\alpha_x(s)^2+\alpha_x(s)A_x(s,s)+O(s^4).
\end{equation}
Differentiation in the first slot, with the second slot held
fixed, gives
\begin{equation}\label{jet:fixed-second}
\begin{split}
D^{(1)}_v b(x,x+s)
={}&-2\alpha_x(s)\alpha_x(v)
   +\alpha_x(s)d\alpha_x(v,s)\\
 &\hspace{8mm}-\alpha_x(v)A_x(s,s)+O(s^3).
\end{split}
\end{equation}
\end{lemma}

\begin{proof}
Write the expansion initially as
$b(x,x+s)=Q_x(s)+C_x(s)+O(s^4)$, where
$Q_x(s)=\alpha_x(s)^2$ and $C_x$ is homogeneous of degree three.
Symmetry, with the same frame in both slots, gives
\[
Q_x(s)+C_x(s)
=Q_{x+s}(-s)+C_{x+s}(-s)+O(s^4).
\]
The two terms on the right are respectively
\[
\alpha_x(s)^2+2\alpha_x(s)A_x(s,s)+O(s^4),
\qquad -C_x(s)+O(s^4).
\]
Thus $C_x(s)=\alpha_x(s)A_x(s,s)$, proving
Equation~\eqref{jet:cubic}.

For the first-slot derivative, $y$ is fixed and hence $s=y-x$
has derivative $-v$. The quadratic term contributes
\[
-2\alpha_x(s)\alpha_x(v)+2\alpha_x(s)A_x(v,s).
\]
Through degree two, the cubic term contributes only through its
arguments $s$, giving
\[
-\alpha_x(v)A_x(s,s)
-\alpha_x(s)\bigl(A_x(v,s)+A_x(s,v)\bigr).
\]
Their sum is Equation~\eqref{jet:fixed-second}, since
$d\alpha_x(v,s)=A_x(v,s)-A_x(s,v)$. Differentiating the holomorphic
remainder with $y$ fixed lowers its order by at most one.

These scalar formulas also keep track of changes of line frame.
If $e'=q e$ and $\rho=q^{-1}$, then
\[
\alpha'=\rho\alpha,
\qquad b'(x,y)=\rho(x)\rho(y)b(x,y).
\]
The cubic coefficient transforms as
\[
C'_x(s)=\rho(x)^2 C_x(s)
+\rho(x)d\rho_x(s)\alpha_x(s)^2
=\alpha'_x(s)(\partial_s\alpha')_x(s).
\]
Moreover the derivative at fixed $y$ transforms by
\[
D^{(1)}_v b'
=\rho(x)\rho(y)D^{(1)}_v b
+d\rho_x(v)\rho(y)b.
\]
In particular, it includes exactly the first-slot frame term.
The cubic is a coordinate expression; the symbol in
Equation~\eqref{jet:second-symbol} is intrinsic.
\end{proof}

\begin{lemma}[Contact Hamiltonians]\label{jet:hamiltonian}
On a complex contact manifold, contact value identifies the
holomorphic vector fields preserving $F$ with $H^0(X,L)$
\cite[Proposition~2.1]{LeBrun1995}.
For $a\in H^0(X,L)$, write $K_a$ for the corresponding field.
In a frame $\theta=e\alpha$, $a=e f$, it is characterized by
\begin{equation}\label{jet:hamiltonian-equations}
\alpha(K_a)=f,
\qquad df(w)+d\alpha(K_a,w)=0\quad(w\in F).
\end{equation}
Its value at a point depends only on the ordinary first jet of $a$.
\end{lemma}

\begin{proof}
Choose a local holomorphic field $R$ with $\alpha(R)=1$ and
write $K_a=fR+Z$ with $Z\in F$. The second equation determines
$Z$ uniquely and holomorphically because $d\alpha|_F$ is
nondegenerate. Cartan's formula identifies that equation with
$(\mathcal L_{K_a}\alpha)|_F=0$, equivalently
$[K_a,F]\subset F$. It also proves that every field preserving
$F$ satisfies the equations for its contact value.

For completeness, under a frame change
$\alpha'=\rho\alpha$, $f'=\rho f$, and for $w\in F$, we have
\[
df'(w)+d\alpha'(K_a,w)
=\rho\bigl(df(w)+d\alpha(K_a,w)\bigr),
\]
since the extra terms are
$d\rho(w)f-d\rho(w)\alpha(K_a)=0$. The first equation transforms
by the same factor. Uniqueness therefore makes the local fields
independent of the chosen frame and lift $R$, and they glue.
The displayed equations use only $f(x)$ and $df_x$ to determine
$K_a(x)$. This also proves naturality under contact automorphisms.
\end{proof}

\begin{theorem}[Symmetric-jet criterion]\label{jet:criterion}
Let $X$ be a connected compact complex contact manifold with
contact quotient $\theta:T_X\to L$, and let $\mathcal I_\Delta$
be the ideal sheaf of the diagonal in $X\times X$.
Suppose that a symmetric section $B\in H^0(X\times X,L\boxtimes L)$
satisfies
\[
\begin{gathered}
B\in H^0(X\times X,\mathcal I_\Delta^2\otimes(L\boxtimes L)),\\
[B]_2=\theta^2
\quad\text{in }H^0(X,\Sym^2\Omega_X^1\otimes L^{\otimes2}).
\end{gathered}
\]
Then the ordinary first-jet evaluation map
\begin{equation}\label{jet:jet-generation}
H^0(X,L)\longrightarrow J_x^1L
\end{equation}
is surjective at every $x\in X$. Global contact vector fields
span $T_xX$ at every point, and the identity component of the
holomorphic contact automorphism group acts transitively.
\end{theorem}

\begin{proof}
Put $d=\dim X$ and $V=H^0(X,L)$, a finite-dimensional vector space
by compactness. Write $B_x$ for the restriction of $B$ to
$\{x\}\times X$, viewed as an element of $L_x\otimes V$.
Point evaluations followed by fiber functionals span $V^*$, since
only the zero section vanishes at every point. Choose a basis
$\varepsilon_1,\ldots,\varepsilon_r$ among these functionals, and let
$s_1,\ldots,s_r$ be its dual basis in $V$. Contracting the second
factor of $B$ by $\varepsilon_j$ gives a global section $t_j$ of $L$
on the first factor. For each $x$, all the $\varepsilon_j$ kill
$B_x-\sum_jt_j(x)\otimes s_j$, so that difference is zero. Hence
$B=\sum_jt_j\otimes s_j$, and in particular
\[
B\in H^0(X,L)\otimes V
=H^0(X,L\otimes(V\otimes\OO_X)).
\]
For each $x$, the vector $B_x\in L_x\otimes V$ is nonzero.
Indeed $B_x=0$ would say that $B(x,y)$ vanishes identically in $y$,
contradicting its nonzero quadratic term $\alpha_x(s)^2$.
Thus $B$ gives a line subbundle and a holomorphic map
\begin{equation}\label{jet:beta}
L^{-1}\hookrightarrow V\otimes\OO_X,
\qquad \beta:X\longrightarrow\PP(V),
\qquad \beta^*\OO_{\PP(V)}(1)\simeq L.
\end{equation}

We claim that $\beta$ is an immersion. In a local frame write
$B=e\otimes w$, with $w$ a holomorphic $V$-valued function.
If $d\beta_x(v)=0$, then $\partial_vw(x)=k w(x)$ for a scalar
$k$, independent of the second variable. Evaluating these global
sections in the frame at $y$ gives
\begin{equation}\label{jet:kernel}
D^{(1)}_v b(x,y)=k b(x,y).
\end{equation}
The linear terms of Equation~\eqref{jet:fixed-second} imply
$\alpha_x(v)=0$, because $\alpha_x$ is a nonzero linear form.
The quadratic terms then give the polynomial identity
\[
\alpha_x(s)d\alpha_x(v,s)=k\alpha_x(s)^2.
\]
Cancelling the nonzero linear polynomial $\alpha_x(s)$ shows
$d\alpha_x(v,\cdot)=k\alpha_x$. Its restriction to $F_x$ is
zero. Since $v\in F_x$, contact nondegeneracy gives $v=0$.

We now pass explicitly from this immersion to actual global
sections. Every $\lambda\in V^*$ gives
\[
t_\lambda=(\id\otimes\lambda)B\in H^0(X,L),
\qquad (t_\lambda)_{\mathrm{loc}}=\lambda(w).
\]
If $v_1,\ldots,v_d$ is a basis of $T_xX$, the $d+1$ vectors
\[
w(x),\ \partial_{v_1}w(x),\ldots,\partial_{v_d}w(x)
\]
are linearly independent: a dependence with nonzero tangent part
would contradict immersion, and a dependence with zero tangent
part contradicts $w(x)\ne0$. Dualizing gives surjectivity of
\[
V^*\longrightarrow L_x\oplus(T_x^*X\otimes L_x),
\qquad
\lambda\longmapsto\bigl(t_\lambda(x),d(t_\lambda)_{\mathrm{loc},x}\bigr),
\]
where the target is the local-frame description of $J_x^1L$.
Recall its intrinsic exact sequence
$0\to\Omega_X^1\otimes L\to J^1L\to L\to0$.
Since every $t_\lambda$ belongs to $H^0(X,L)$, this proves
Equation~\eqref{jet:jet-generation}. No nondegeneracy of $B$ as
a tensor in $V\otimes V$ is required.

Given $u\in T_xX$, prescribe a first jet $j_x^1a$ whose scalar
representative satisfies
\[
f(x)=\alpha_x(u),
\qquad df_x|_{F_x}=-d\alpha_x(u,\cdot)|_{F_x}.
\]
Extend the prescribed covector arbitrarily from $F_x$ to $T_xX$.
First-jet surjectivity realizes this jet by a global section $a$.
Lemma~\ref{jet:hamiltonian} and uniqueness in
Equation~\eqref{jet:hamiltonian-equations} imply $K_a(x)=u$.
Thus global contact fields span every tangent space.

A global holomorphic vector field on a compact complex manifold
has a global flow for sufficiently small complex time. Compactness
allows a common time disk for all initial points, and the
inverse is supplied by negative time. For a contact field these
automorphisms preserve $F$ and form a path through the identity.
At any point, compositions of such flows for a spanning set of
fields have surjective differential at zero. Every orbit of the
identity component therefore contains a neighborhood of each of
its points. All orbits are open, and connectedness of $X$ implies
that there is only one.
\end{proof}

\subsection{The algebraic action}

\begin{corollary}\label{jet:transitive}
For every smooth contact Fano manifold in Theorem~\ref{thm:main},
the identity component $G=\Aut(X,F)^0$ is a connected linear
algebraic group acting transitively on $X$.
\end{corollary}

\begin{proof}
In the primitive Picard case, Proposition~\ref{jet:diagonal} and
Theorem~\ref{jet:criterion} give holomorphic contact transitivity.
Proposition~\ref{red:alternatives} supplies a connected transitive
contact group in both remaining Picard alternatives.

Choose $m>0$ such that $-mK_X$ is very ample. Its canonical
functoriality embeds $\Aut(X)$ as the closed stabilizer of $X$ in
the projective linear group of this anticanonical embedding.
Preserving $F$ is also a closed condition: under the induced
action on $\PP(T_X)$ it is the condition of stabilizing the closed
subvariety $\PP(F)$. Thus $\Aut(X,F)$ is linear algebraic.
Holomorphic automorphisms of projective $X$ are algebraic, and
the contact flows in the proof of Theorem~\ref{jet:criterion}
belong to its identity component: their action on the
anticanonical embedding varies continuously from the identity.
Consequently $G$ acts transitively in the primitive case as well.
\end{proof}
